
%\needspace{4\baselineskip}
\color{uniblue}{\subsection{Логика отключения/включения (ЛОВ)}}
\color{black}

\begin{enumerate}
%%%%%%%%%%%%%%%%%%%%%%%%%%%%%%%%%%%%%%%%%%%%%%%%%%%%%%%%%%%  ОБЩИЕ СВЕДЕНИЯ %%%%%%%%%%%%%%%%%%%%%%%%%%%%%%%%%%%%%%%%%%%%%%%%

\item Общие сведения

\begin{enumerate}

\item
Функция логики отключения/включения выключателя предназначена для подачи аварийных команд управления на выключатель при действии устройств РЗА. В устройстве логика представлена функциональным блоком «ЛОВ».

\item
Функциональный блок состоит из следующих функций:
\begin{itemize}
\item логика отключения от релейной защиты (ЛО РЗ); 
\item логика отключения ВНР (ЛО ВНР);
\item Логика включения ВНР (ЛВ ВНР);
\item логика запрета АПВ (ЗАПВ);
\item Логика запрета АВР (ЗАВР ЛО).
\end{itemize}

\item
Функциональный блок может быть полностью выведен из работы оперативно путем активации сигнала «ОВ ЛОВ».

\item
В таблице \ref{lorz:tbl1} приведены параметры, необходимые для настройки ЛОВ.

\end{enumerate}

%%%%%%%%%%%%%%%%%%%%%%%%%%%%%%%%%%%%%%%%%%%%%%%%%%%%%%%%% Логика отключения от релейной защиты (ЛО РЗ) %%%%%%%%%%%%%%%%%%%%%%%%%%%%%%%%%%%%%%%%%%%%%%%%

\item Логика отключения от релейной защиты (ЛО РЗ)

\begin{enumerate}

\item
ЛО РЗ формирует сигналы аварийного отключения («ЛО РЗ: Отключение», «ЛО РЗ: Отключить») при срабатывании функций релейной защиты в составе устройства, а также от четырех назначаемых сигналов внешнего отключения («Откл. внешнее 1-4»). В функциональном блоке «ЛОВ» логика отключения от релейной защиты представлена функцией «ЛО РЗ».

\item
Сигнал «ЛО РЗ: Отключение» является суммарным сигналом отключения от отдельных функций и существует пока активен хотя бы один сигнал срабатывания защитных функций.

\item
Сигнал «ЛО РЗ: Отключить» воздействует на выходные реле для отключения выключателя, длительность сигнала задается уставкой Т1.

\item
Ввод функции в работу на этапе параметрирования устройства осуществляется программным переключателем SGF1 «Ввод\_функции». Информация о введенной функции отображается выходным сигналом «ЛО РЗ: Ввод».

\item
Функция может быть выведена из работы оперативно путем активации сигнала «ОВ ЛО РЗ», что характеризуется активным состоянием сигнала «ЛО РЗ: Оперативный вывод» на выходе функции.

\item
Алгоритм работы функции выполнен в соответствии с рисунком \ref{lorz:img1}.

\vspace{1mm}
\begin{figure}[H]
\centering
\includegraphics[width=0.95\textwidth,height=0.95\textheight,keepaspectratio]{\fbpath/311. ЛОВ ВВ/_latex/img/lorz.pdf}
\captionof{figure}{Функциональная схема ЛО РЗ}\label{lorz:img1}
\end{figure}


\end{enumerate}

%%%%%%%%%%%%%%%%%%%%%%%%%%%%%%%%%%%%%%%%%%%%%%%%%%%%%%%%%  Логика отключения ВНР (ЛО ВНР)  %%%%%%%%%%%%%%%%%%%%%%%%%%%%%%%%%%%%%%%%%%%%%%%%

\item  Логика отключения ВНР (ЛО ВНР) 

\begin{enumerate}

\item\sloppy
ЛО ВНР формирует сигналы аварийного отключения («ЛО ВНР: Отключение», «ЛО ВНР: Отключить») при срабатывании внешней функции ВНР «Отключение от ВНР». В функциональном блоке «ЛОВ» логика отключения ВНР представлена функцией «ЛО ВНР». 

\item
Сигнал «ЛО ВНР: Отключить» воздействует на выходные реле для отключения выключателя, длительность сигнала задается уставкой Т1.

\item
Ввод функции в работу на этапе параметрирования устройства осуществляется программным переключателем SGF1 «Ввод\_функции». Информация о введенной функции отображается выходным сигналом «ЛО ВНР: Ввод».

\item
Функция может быть выведена из работы оперативно путем активации сигнала «ОВ ВНР», что характеризуется активным состоянием сигнала «ЛО ВНР: Оперативный вывод» на выходе функции.

\item
Алгоритм работы функции выполнен в соответствии с рисунком \ref{lovnr:img1}.

\vspace{1mm}
\begin{figure}[H]
\centering

\includegraphics[width=0.9\textwidth,height=0.9\textheight,keepaspectratio]{\fbpath/311. ЛОВ ВВ/_latex/img/lovnr.pdf}
\captionof{figure}{Функциональная схема ЛО ВНР}\label{lovnr:img1}
\end{figure}


\end{enumerate}
%%%%%%%%%%%%%%%%%%%%%%%%%%%%%%%%%%%%%%%%%%%%%%%%%%%%%%%%% КОНЕЦ Логика отключения ВНР (ЛО ВНР) %%%%%%%%%%%%%%%%%%%%%%%%%%%%%%%%%%%%%%%%%%%%%%%%

%%%%%%%%%%%%%%%%%%%%%%%%%%%%%%%%%%%%%%%%%%%%%%%%%%%%%%%%%  Логика включения ВНР (ЛВ ВНР)  %%%%%%%%%%%%%%%%%%%%%%%%%%%%%%%%%%%%%%%%%%%%%%%%

\item Логика включения ВНР (ЛВ ВНР) 

\begin{enumerate}

\item
ЛВ ВНР формирует сигнал включения («ЛВ ВНР: Включить») при срабатывании внешнего ВНР на включение присоединения. В функциональном блоке «ЛОВ» логика включения ВНР представлена функцией «ЛВ ВНР». 

\item
Функция предусматривает включение выключателя с контролем синхронизма. Тип пуска определяется уставкой SGF2 «Контр\_синхр».

\item
Включение выключателя с контролем синхронизма осуществляется при наличии входного сигнала разрешения автоматического включения выключателя из логики КС «Разреш.авт.вкл. В» и поступлении на дискретный  вход  устройства  сигнала  от  секционного  выключателя  «Включение  от  ВНР»  на  время  не регулируемого таймера Т2 равное 5 мс при отсутствии сигналов для сброса триггера. 

\item
Включение  выключателя  без  контроля  синхронизма  осуществляется  только  при  поступлении  на дискретный  вход  устройства  сигнала  от  секционного  выключателя  «Включение  от  ВНР»  на  время  не регулируемого таймера Т2 равное 5 мс при отсутствии сигналов для сброса триггера. 

\item
При алгоритме включения выключателя с КС сброс триггера будет осуществляться при задержке включения выключателя по истечении выдержки времени таймера Т1. Прохождение команды включения выключателя  «ЛВ  ВНР:  Включить» так  же сбрасывает  триггер.  Дополнительными  условиями  сброса триггера является наличие/отсутствие сигналов из логики ЗАВР и КСВ.

\item
Ввод функции в работу на этапе параметрирования устройства осуществляется программным переключателем SGF1 «Ввод\_функции». Информация о введенной функции отображается выходным сигналом «ЛВ ВНР: Ввод».

\item
Функция может быть выведена из работы оперативно путем активации сигнала «ОВ ЛВ ВНР», что характеризуется активным состоянием сигнала «ЛВ ВНР: Оперативный вывод» на выходе функции.

\item
Алгоритм работы функции выполнен в соответствии с рисунком \ref{lvvnr:img1}. 

\vspace{1mm}
\begin{figure}[H]
\centering
\includegraphics[width=1\textwidth,height=1\textheight,keepaspectratio]{\fbpath/311. ЛОВ ВВ/_latex/img/lvvnr.pdf}
\captionof{figure}{Функциональная схема ЛВ ВНР}\label{lvvnr:img1}
\end{figure}

\end{enumerate}
%%%%%%%%%%%%%%%%%%%%%%%%%%%%%%%%%%%%%%%%%%%%%%%%%%%%%%%%% КОНЕЦ  Логика включения ВНР (ЛВ ВНР)  %%%%%%%%%%%%%%%%%%%%%%%%%%%%%%%%%%%%%%%%%%%%%%%%

%%%%%%%%%%%%%%%%%%%%%%%%%%%%%%%%%%%%%%%%%%%%%%%%%%%%%%%%% Запрет АПВ %%%%%%%%%%%%%%%%%%%%%%%%%%%%%%%%%%%%%%%%%%%%%%%%

\item Логика запрета АПВ (ЗАПВ)

\begin{enumerate}

\item
Запрет АПВ может быть необходим в определённых ситуациях для предотвращения повторных срабатываний и возможных повреждений оборудования. В функциональном блоке «ЛОВ» логика запрета АПВ представлена функцией «ЗАПВ». 

\item
Сигнал «ЗАПВ: Запрет АПВ» формируется, если присутствует сигнал срабатывания защит на любом из входов, которые заведены под запрет АПВ с помощью выбора уставки программного переключателя. Контроль всех входных сигналов, выполняется соответствующими уставками (SGF1 \ldots SGF23).

\item
Ввод функции в работу осуществляется при введеной функции АПВ в работу («АПВ: Ввод»). Информация о введенной функции отображается выходным сигналом «ЗАПВ: Ввод».

\item
Функция может быть выведена из работы оперативно путем активации сигнала «ОВ ЛОВ», что характеризуется активным состоянием сигнала «ЗАПВ: Оперативный вывод» на выходе функции.

\item
Алгоритм работы функции выполнен в соответствии с рисунком \ref{zapv:img1}. 

\vspace{1mm}
\begin{figure}[H]
\centering
\includegraphics[width=0.8\textwidth,height=0.8\textheight,keepaspectratio]{\fbpath/311. ЛОВ ВВ/_latex/img/zapv.pdf}
\captionof{figure}{Функциональная схема логики ЗАПВ}\label{zapv:img1}
\end{figure}

\end{enumerate}
%%%%%%%%%%%%%%%%%%%%%%%%%%%%%%%%%%%%%%%%%%%%%%%%%%%%%%%%% КОНЕЦ Запрет АПВ %%%%%%%%%%%%%%%%%%%%%%%%%%%%%%%%%%%%%%%%%%%%%%%%

%%%%%%%%%%%%%%%%%%%%%%%%%%%%%%%%%%%%%%%%%%%%%%%%%%%%%%%%% Запрет АВР %%%%%%%%%%%%%%%%%%%%%%%%%%%%%%%%%%%%%%%%%%%%%%%%

\item Логика запрета АВР (ЗАВР ЛО)

\begin{enumerate}

\item
Запрет АВР может быть необходим в определённых ситуациях для предотвращения повторных срабатываний и возможных повреждений оборудования. В функциональном блоке «ЛОВ» логика запрета АВР представлена функцией «ЗАВР ЛО».

\item
Сигнал «ЗАВР: Запрет АВР» формируется, если присутствует сигнал срабатывания защит на любом из входов, которые заведены под запрет АВР с помощью выбора уставки программного переключателя. Контроль всех входных сигналов внешнего отключения и сигналов «КП / УВ: Отключить», «Внешний запрет АВР» выполняется соответствующими уставками (SGF2 \ldots SGF7).

\item
Ввод функции в работу осуществляется при введенной функции АВР в работу («АВР: Ввод»). Информация о введенной функции отображается выходным сигналом «ЛОВ / ЗАВР: Ввод».

\item
Функция может быть выведена из работы оперативно путем активации сигнала «ОВ ЛОВ», что характеризуется активным состоянием сигнала «ЛОВ / ЗАВР: Оперативный вывод» на выходе функции.

\item
Алгоритм работы функции выполнен в соответствии с рисунком \ref{zavr:img1}.

\vspace{3mm}
\begin{figure}[H]
\centering

\includegraphics[width=0.8\textwidth,height=0.8\textheight,keepaspectratio]{\fbpath/311. ЛОВ ВВ/_latex/img/zavr.pdf}
\captionof{figure}{Функциональная схема логики ЗАВР}\label{zavr:img1}
\end{figure}

\end{enumerate}
%%%%%%%%%%%%%%%%%%%%%%%%%%%%%%%%%%%%%%%%%%%%%%%%%%%%%%%%% КОНЕЦ Запрет АПВ %%%%%%%%%%%%%%%%%%%%%%%%%%%%%%%%%%%%%%%%%%%%%%%%


\end{enumerate}
